\workbookchapter{混合专家模型}

\begin{exercises}
\begin{exercise} 对 $E_r=4,k=2$、分数 $(\log0.4,\log0.3,\log0.2,\log0.1)$，计算选中内归一化和保留全专家概率的两种权重。设已选专家输出为 $(1,0)$ 与 $(0,2)$，比较最终输出。

\begin{solution}
选中专家1、2。选中内权重为 $(\frac{4}{7},\frac{3}{7})$，输出 $(\frac{4}{7},\frac{6}{7})$；保留全专家概率时为(.4,.3)，输出(.4,.6)。后者权重总和.7，缺少的.3没有被重分配，所以两种输出不同；选择集合相同不代表函数相同。
\end{solution}
\end{exercise}

\begin{exercise} 从式\tbobject{eq:rev11-router-grad}证明：固定选中集合且 $k=1$ 时，选中内归一化的路由任务梯度为零。再解释保留全专家概率为何不必为零。

\begin{solution}
固定集合、选中归一化时 $y=\sum_iw_iF_i$，对选中logit的梯度为 $w_jg^{\mathsf T}(F_j-y)$。$k=1$时 $w_j=1,y=F_j$，因此为零。若保留全专家Softmax概率，$y=p_jF_j$，$p_j$ 仍受所有logit影响，故通常有非零梯度。选择边界本身不可微，局部零梯度不表示改变排名永远不改变输出。
\end{solution}
\end{exercise}

\begin{exercise}\optional 对某个词元，两个专家输出相同。分析选中内归一化时路由权重的任务梯度，并说明这是否意味着专家参数梯度也为零。

\begin{solution}
两个选中专家均输出v时，归一化混合始终为v，$F_j-y=0$，路由任务梯度为零。但专家输出的上游梯度分别是 $w_jg$，专家参数通过其雅可比仍可得到非零梯度。路由选择无差别不等于专家停止学习；辅助均衡也可能继续更新路由。
\end{solution}
\end{exercise}

\begin{exercise} 构造 $E_rC\ge Nk$ 但仍存在容量溢出的例子。分别计算路径拒绝率、受影响词元比例和完全失去专家处理的比例。

\begin{solution}
取$N=2$、$k=1$、$E=2$、每专家容量$C=1$，则 $EC=Nk=2$。若两个词元都选择专家1，按词元编号接收一个、拒绝另一个。路径拒绝率$\frac{1}{2}$，受影响词元率$\frac{1}{2}$，完全失去专家处理率$\frac{1}{2}$。专家2的空闲槽位不会自动把已选路径变为合法替代；重新路由是另一策略。
\end{solution}
\end{exercise}

\begin{exercise}\optional 在本章数值例题中，把专家2接收规则改为按该专家的路由权重降序选择，并列时按词元编号升序接收。确定哪条路径被拒绝，再比较输出变化。说明规则不同为何不能只当成排序实现细节。

\begin{solution}
专家2的四个权重为 $\frac{2}{9},\frac{1}{3},\frac{7}{18},\frac{2}{9}$。按降序为词元3、2，然后词元1与4并列；题设补充并列按词元编号升序，故仍拒绝词元4，输出与正文的同序接收例相同。若并列优先较大编号，则拒绝词元1：不重归一化时 $y_1=(\frac{7}{9})x_1$，其余恢复原无拒绝输出。该例中的是否改变依赖明确的并列规则；不能凭方法名称声称一定改善。
\end{solution}
\end{exercise}

\begin{exercise} 推导式\tbobject{eq:rev11-balance-grad}。如果辅助统计覆盖两个大小不同的本地批次，写出正确的全局平均概率，解释局部均值直接平均的问题。

\begin{solution}
设 $L=E\sum_if_i\bar p_i$，$\bar p_i=N^{-1}\sum_tp_{ti}$，并对离散统计f停止梯度。由Softmax导数得 $\frac{\partial L}{\partial z_{tj}}=(\frac{E}{N})p_{tj}(f_j-\sum_if_ip_{ti})$。两个批次大小 $N_1,N_2$ 时，正确全局概率为 $\frac{N_1\bar p^{(1)}+N_2\bar p^{(2)}}{N_1+N_2}$；直接各占一半将小批次词元过度加权，路由计数也须在同一统计范围聚合。
\end{solution}
\end{exercise}

\begin{exercise} 令 $d=512,d_f=2048,E_r=16,k=2$，忽略偏置，计算门控专家的总参数与每词元激活参数。加入两个始终执行的同尺寸共享专家后，两个计数分别如何变化？

\begin{solution}
一个无偏置门控专家含三矩阵，共 $3dd_f=3145728$ 参数。16个路由专家总计50331648，每词元选2个执行6291456。新增2个同尺寸共享专家后，总量56623104，激活量12582912。这里只数专家矩阵；若计入路由器还需加 $dE=8192$，注意力、归一化及输出头另计。
\end{solution}
\end{exercise}

\begin{exercise}\optional 保持每词元专家矩阵计算近似不变，将专家中间宽度减半、专家数与选择数加倍。哪些成本不能随专家矩阵算术量保持不变？

\begin{solution}
$d_f$减半、E和k加倍保持 $3kdd_f$ 及专家总矩阵参数近似不变，但路由logit数量翻倍，选择/分发元数据和路径数增多。每专家矩阵更窄、每路径批次更小，内核利用率可能下降；通信消息数量、负载尾部与合并开销也变化。需同硬件同token预算测量，不能由乘加次数守恒推出时延守恒。
\end{solution}
\end{exercise}

\begin{exercise}\optional 对 Soft MoE 分发与合并矩阵，分别验证列和与行和。举例说明若先混合全序列词元再调用专家，为什么早期输出可能读取未来输入。

\begin{solution}
若分发权重 $D_{ts}=\frac{\exp z_{ts}}{\sum_u\exp z_{us}}$，则每个槽位列和为1；合并权重 $C_{ts}=\frac{\exp z_{ts}}{\sum_r\exp z_{tr}}$，则每个词元行和为1。两词元一个槽位、权重各$\frac{1}{2}$时，槽输入 $\frac{x_1+x_2}{2}$，若专家为恒等，则第一个输出直接含x2。仅在最终合并处加因果mask不能消除已经混入的未来信息。
\end{solution}
\end{exercise}

\begin{exercise} 某模型在训练大批次中专家利用率较高，在低并发解码中却很慢。用矩阵形状、权重读取、通信延迟和路由负载解释至少三种可能机制。

\begin{solution}
低并发使每专家输入行数很少，矩阵乘法退化为低利用率的小矩阵；大量专家权重可能仍需从显存读取，计算量低但带宽占主导。跨设备每轮小消息暴露固定通信延迟，路由偏斜又使少数专家成为串行尾部。应记录逐专家批量、读取字节、通信耗时和请求尾延迟，再决定合批、放置或卸载策略。
\end{solution}
\end{exercise}

\begin{exercise} 某报告只给出平均专家负载和激活参数量，便断言单卡内存与吞吐满足上线需求。说明尚缺哪些可支持结论的结构信息和运行证据。

\begin{solution}
需要逐层专家数/容量/共享布局、总驻留权重精度、路由和KV状态、峰值工作区、设备分片及卸载方案；还需请求长度并发、批次分布、溢出与重试、内核和通信数据。以实际候选模型在目标硬件测吞吐、首词延迟和尾延迟，并验证任务质量。平均负载会掩盖热点，激活参数也不等于显存。
\end{solution}
\end{exercise}

\end{exercises}
